% ==============================================================================
% PDF-1 / Section 05 — STOP Regions (Skeleton)
% No new primitives. Diagnostic-only. Container structure.
% ==============================================================================

\section{STOP Regions}
\label{sec:stop-regions}

\subsection{Conceptual role of STOP regions}
STOP regions represent formally defined boundaries in the enterprise phase space
beyond which diagnostic execution must halt.
They are first-class diagnostic outcomes, not errors and not recommendations.

STOP regions encode structural impossibility, loss of admissibility, or violation
of invariant constraints already specified in the KEA and ETS layers.

\subsection{STOP as a diagnostic outcome}
Reaching a STOP region signifies that further diagnostic refinement is either
undefined or non-falsifiable under the given specification.
This outcome is explicit, terminal, and non-negotiable.

STOP does not imply failure, blame, urgency, or the necessity of action.

\subsection{Boundary characterization (container)}
STOP regions are characterized by:
\begin{itemize}
  \item violation of admissibility constraints,
  \item collapse of observability assumptions,
  \item incompatibility with invariant stress conditions.
\end{itemize}

No new boundary criteria are introduced here.
All references are strictly inherited from immutable sources.

\subsection{Detection and signaling}
This section reserves the structural location for detecting proximity to, or entry
into, STOP regions during diagnostic execution.
Signaling of STOP is treated as a formal result, not a warning or alert.

\subsection{Explicit exclusions}
\begin{itemize}
  \item No recovery strategies.
  \item No mitigation or workaround logic.
  \item No escalation or governance prescriptions.
  \item No optimization away from STOP.
\end{itemize}
